%%%%%%%%%%%%%%%%%%%%%%%%%%%%%%%%%%%% Section 2: Related Work %%%%%%%%%%%%%%%%%%%%%%%%%%%%%%%%%%%%
\section{Related Work}
\label{sec:Related_Work}

The Rapidly-exploring Random Tree (RRT) algorithm \cite{lavalle1998rapidly} is a sampling-based method for path planning, designed to efficiently explore high-dimensional configuration spaces. It incrementally builds a tree by randomly sampling points in the space and connecting them to the nearest point in the existing tree, ensuring rapid exploration. RRT is particularly effective for finding feasible paths in complex, obstacle-filled environments, making it a widely used approach in robotics and autonomous navigation.

Significant algorithmic refinements have enhanced RRT's capabilities over the years. For instance, the RRT* algorithm \cite{karaman2011sampling} introduced asymptotic optimality through node-rewiring mechanisms by minimizing a predefined cost function (e.g., path length), while Informed RRT* \cite{gammell2014informed} developed advanced sampling strategies to guide exploration near optimal solution regions. On the other hand, RRT-Connect \cite{kuffner2000rrt} grows two trees bidirectionally from both the start and goal configurations. This approach, combined with a greedy heuristic, allows faster exploration of the configuration space and quicker convergence to a solution compared to the standard RRT.

To tackle issues related to random sampling and low path efficiency in RRT, an improved approach incorporating the Artificial Potential Field (APF) method was recently proposed in \cite{huang2021path}. This method introduces a probability value during the expansion step of the random tree in the basic RRT algorithm, enhancing convergence speed toward the target node. Similarly, the authors in \cite{app14010025} proposed an RRT variant that utilizes adjustable probability and sampling area strategies to quickly find a feasible path. This planner is then combined with an optimizer that uses the Dijkstra algorithm to prune and improve the initial path. More recently, the work in \cite{liu2024rrt} combined RRT with model predictive control (MPC) utilizing control barrier functions (CBF) to enforce safety-critical constraints. Additionally, recent advancements in learning-based RRT methods have shown promising results in improving path-planning efficiency. The Neural Informed RRT* approach \cite{huang2024neural} utilizes a neural network to learn the topology of the free space and infer states close to the optimal path, thereby guiding the search toward more promising regions, whereas Neural RRT* \cite{wang2020neural} employs convolutional neural networks to predict a probabilistic heatmap of states for guiding exploration.

Finally, the emergence of Large Language Models (LLMs) \cite{touvron2023llama, brown2020language} has revolutionized artificial intelligence, enabling advanced reasoning and knowledge-driven applications in autonomous navigation. Recent research has explored the potential of LLMs to enhance navigation tasks by combining their reasoning capabilities with domain-specific methods. For instance, the work in \cite{zhou2024navgpt} demonstrated that LLMs can perform high-level planning tasks for navigation, including identifying landmarks from observed scenes, tracking navigation progress, and correcting course. More closely related to our work is the approach in \cite{meng2024llm}, where the authors integrated LLMs with the A* path-planning algorithm \cite{Hart1968}, achieving enhanced pathfinding efficiency in terms of time and space complexity.

In summary, RRT-based path-planning approaches are extensively utilized for their ability to explore state spaces efficiently and effectively. However, they often require meticulous parameter tuning and tend to exhibit slow convergence. While these methods can find optimal solutions, they frequently incur significant computational overhead, with high memory and time demands, particularly when searching for the best path. This limitation is especially critical in applications such as autonomous vehicles, where rapid identification of efficient paths is essential due to constraints like limited power or fuel resources. Motivated by these challenges, this work proposes a novel path-planning framework that integrates the reasoning capabilities of multimodal large language models with RRT-based path search to enhance the efficiency of path generation.


